\usetikzlibrary{shapes.geometric,arrows.meta,calc,babel,backgrounds,decorations.pathreplacing}
\usepgfplotslibrary{statistics}

\tikzset{
    >={Triangle[angle=60:4]},
        help help lines/.style={lightgray!50,very thin},
        help lines/.style={gray!50,thin},
    barbarrow/.style={
        >={Straight Barb[length=5pt,width=5pt]}
        },
    block/.style={
        rectangle with rounded corners,
        minimum width=2.5cm,
        minimum height=1.1cm,
        inner sep=0pt,
        align=center,
        draw=black,
        fill=white
    },
    db/.style={
        cylinder,
        thick,
        draw = black,
        aspect = .125,
        % minimum height = 1.2cm,
        inner sep=5pt,
        minimum width = 1.5cm,
        shape border rotate = 90
    },
    umlClass/.style={
        fill=white,
        draw=black,
        align=center,
        inner sep=5pt
    },
    M/.style={
        rectangle,
        rounded corners,
        minimum width=2cm,
        inner sep=3pt,
        align=center,
        draw=tubaf-black,
        semithick,
    },
    MA/.style={
        M,
        ellipse,
    },
    MP/.style={
        M,
        rectangle,
        rounded corners = false
    },
    G0/.style={draw=tubaf-primary},
    E0/.style={G0, fill=tubaf-primary, thick, rounded corners},
    Ei/.style={thick, dashed, rounded corners},
    Relation/.style={thick, dashed, tubaf-black},
    Tree/.style={Ei, draw=gray},
    % TreeNamespace/.style={Ei, draw=tubaf-accent-4},
    TreeNamespace/.style={Ei, draw=tubaf-environment!35!tubaf-accent-4},
    TreeExecutor/.style={Ei, draw=tubaf-accent-24},
    TreeContainer/.style={Ei, draw=tubaf-accent-13},
    TreeProcess/.style={Ei, draw=tubaf-energy},
    rectangle with rounded corners north west/.initial=4pt,
    rectangle with rounded corners south west/.initial=4pt,
    rectangle with rounded corners north east/.initial=4pt,
    rectangle with rounded corners south east/.initial=4pt,
}

\pgfplotsset{
    colormap={colormap-cpu}{
        color(0pt)=(tubaf-primary-dark);
        color(20pt)=(tubaf-energy);
        color(40pt)=(tubaf-accent-12);
        color(60pt)=(tubaf-accent-15);
        color(80pt)=(tubaf-accent-17);
        color(100pt)=(tubaf-white);
    },
    colormap={colormap-ram}{
        color(0pt)=(tubaf-primary-dark);
        color(25pt)=(tubaf-material);
        color(50pt)=(tubaf-accent-29);
        color(75pt)=(tubaf-accent-25);
        color(100pt)=(tubaf-white);
    },
    colormap={colormap-net}{
        color(0pt)=(tubaf-primary-dark);
        color(20pt)=(tubaf-environment);
        color(50pt)=(tubaf-accent-3);
        color(75pt)=(tubaf-accent-5);
        color(100pt)=(tubaf-white);
    },
    colormap={colormap-latencies}{
        color(0pt)=(tubaf-primary-dark);
        color(.05pt)=(tubaf-white);
    },
    colormap={colormap-disk}{
        color(0pt)=(tubaf-primary-dark);
        color(25pt)=(tubaf-accent-11);
        color(50pt)=(tubaf-accent-13);
        color(75pt)=(tubaf-accent-16);
        color(100pt)=(tubaf-white);
    },
    cpu/.style={
        xlabel={Frequency [Hz]},
        ylabel={Number of Pub-Sub pairs},
        tick align=center,
        colormap name=colormap-cpu,
        colorbar,
        colorbar sampled,
        colorbar style={
            samples=512,
            ylabel=CPU usage [\%],
            yticklabel style={align=right},
        },
        xtick pos=bottom,
        ytick pos=left,
        % point meta min=0,
        % point meta max=100,
    },
    ram/.style={
        xlabel={Frequency [Hz]},
        ylabel={Number of Pub-Sub pairs},
        tick align=center,
        colormap name=colormap-ram,
        colorbar,
        colorbar sampled,
        colorbar style={
            samples=512,
            ylabel=RAM usage [MB],
            yticklabel style={align=right},
        },
        xtick pos=bottom,
        ytick pos=left,
        % point meta min=0,
        % point meta max=1500,
    },
    net/.style={
        xlabel={Frequency [Hz]},
        ylabel={Number of Pub-Sub pairs},
        tick align=center,
        colormap name=colormap-net,
        colorbar,
        colorbar sampled,
        colorbar style={
            samples=512,
            ylabel=Network usage [kB/s],
            yticklabel style={align=right},
        },
        xtick pos=bottom,
        ytick pos=left,
        % point meta min=0,
        % point meta max=1500,
    },
    disk/.style={
        xlabel={Frequency [Hz]},
        ylabel={Number of Pub-Sub pairs},
        tick align=center,
        colormap name=colormap-disk,
        colorbar,
        colorbar sampled,
        colorbar style={
            samples=512,
            ylabel=Disk Write [kB/s],
            yticklabel style={align=right},
        },
        xtick pos=bottom,
        ytick pos=left,
        % point meta min=0,
        % point meta max=1500,
    },
    lat/.style={
        xlabel={Frequency [Hz]},
        ylabel={Number of Pub-Sub pairs},
        colormap name=colormap-latencies,
        colorbar,
        colorbar sampled,
        colorbar style={
            samples=512,
            ylabel=Overhead [$\upmu \text{s}$],
            yticklabel style={
                align=right,
            }
        }
    }
}


\newcommand{\activeMember}[3]{
    \node [circle, ultra thick, draw=tubaf-primary, fill=cyan] (#1) at #2 {\phantom{foooooo}};
    \draw (#1) node {\footnotesize#3};
}
\newcommand{\activeCausalMember}[3]{
    \node [circle, ultra thick, draw=tubaf-primary, fill=cyan] (#1) at #2 {\phantom{foooooo}};
    \draw (#1) node {\footnotesize\ul{#3}};
}
\newcommand{\passiveMember}[3]{
    \node [circle, very thick, draw=tubaf-accent-20, fill=tubaf-accent-23] (#1) at #2 {\phantom{foooooo}};
    \draw (#1) node {\footnotesize#3};
}
\newcommand{\passiveCausalMember}[3]{
    \node [circle, very thick, draw=tubaf-accent-20, fill=tubaf-accent-23] (#1) at #2 {\phantom{foooooo}};
    \draw (#1) node {\footnotesize\ul{#3}};
}
\newcommand{\CommunicationTimer}[2]{
    \draw [G0, thick] (#1.#2+10) -- ++(#2+10:1.1) arc (#2+100:#2:.352);
    \draw [G0, thick, <-] (#1.#2-10) -- ++(#2-10:1.1) arc (#2-100:#2:.352) coordinate (tmp);
    \draw (#1) ++(#2:1.8) coordinate (tmp);
    \node at (tmp) [rotate=#2-180] {\scriptsize timer};
}
\newcommand{\CommunicationTimerNoText}[2]{
    \draw [G0, thick] (#1.#2+10) -- ++(#2+10:1.1) arc (#2+100:#2:.352);
    \draw [G0, thick, <-] (#1.#2-10) -- ++(#2-10:1.1) arc (#2-100:#2:.352) coordinate (tmp);
}
\newcommand{\curvedEdge}[4]{
    \draw[->, bend left=#4] (#1) to (#2);
    \ifthenelse{\isempty{#3}}
        {}
        {\path[bend left=#4] (#1) to node[midway, fill=white, inner sep=1pt] {\scriptsize #3} (#2);}
}
\newcommand{\curvedEdgeWithProperties}[5]{
    \draw[->, bend left=#4, #5] (#1) to (#2);
    \ifthenelse{\isempty{#3}}
        {}
        {\path[bend left=#4] (#1) to node[midway, fill=white, inner sep=1pt] {\scriptsize #3} (#2);}
}


\newcommand{\whitebigodot}[2]{
    \node [circle, thick, draw=#1, fill=white, scale=.5] (a) at #2 {\phantom{$\cdot$}};
    \node (a) at #2 {\color{#1}$\cdot$};
}
\newcommand{\whitebigotimes}[2]{
    \node [circle, thick, draw=#1, fill=white, scale=.5] (a) at #2 {\phantom{$\cdot$}};
    \draw [thick, #1] (a.south west) -- (a.north east);
    \draw [thick, #1] (a.north west) -- (a.south east);
}

\newcommand{\trianglearrow}{
  \mathrel{
    \tikz[baseline=-0.5ex]
      \draw[->, line width=0.4pt] (0,0) -- (1em,0);
  }
}


\makeatletter
\pgfdeclareshape{rectangle with rounded corners}{
\inheritanchorborder[from=rectangle]
\savedmacro{\neoffset}{
    \pgfkeysgetvalue{/tikz/rectangle with rounded corners north east}{\pgf@rectc}
    \let\neoffset\pgf@rectc
}
\savedmacro{\nwoffset}{
    \pgfkeysgetvalue{/tikz/rectangle with rounded corners north west}{\pgf@rectc}
    \let\nwoffset\pgf@rectc
}
\savedmacro{\seoffset}{
    \pgfkeysgetvalue{/tikz/rectangle with rounded corners south east}{\pgf@rectc}
    \let\seoffset\pgf@rectc
}
\savedmacro{\swoffset}{
    \pgfkeysgetvalue{/tikz/rectangle with rounded corners south west}{\pgf@rectc}
    \let\swoffset\pgf@rectc
}
\savedanchor{\north}{
    \pgf@y=.5\ht\pgfnodeparttextbox
    \pgf@x=0pt
    \setlength{\pgf@ya}{\pgfshapeminheight}
    \ifdim\pgf@y<.5\pgf@ya
    \pgf@y=.5\pgf@ya
    \fi
}
\savedanchor{\south}{
    \pgf@y=-.5\ht\pgfnodeparttextbox
    \pgf@x=0pt
    \setlength{\pgf@ya}{\pgfshapeminheight}
    \ifdim\pgf@y>-.5\pgf@ya
    \pgf@y=-.5\pgf@ya
    \fi
}
\savedanchor{\east}{
    \pgf@y=0pt
    \pgf@x=.5\wd\pgfnodeparttextbox
    \addtolength{\pgf@x}{2ex}
    \setlength{\pgf@xa}{\pgfshapeminwidth}
    \ifdim\pgf@x<.5\pgf@xa
    \pgf@x=.5\pgf@xa
    \fi
}
\savedanchor{\west}{
    \pgf@y=0pt
    \pgf@x=-.5\wd\pgfnodeparttextbox
    \addtolength{\pgf@x}{-2ex}
    \setlength{\pgf@xa}{\pgfshapeminwidth}
    \ifdim\pgf@x>-.5\pgf@xa
    \pgf@x=-.5\pgf@xa
    \fi
}
\savedanchor{\northeast}{
    \pgf@y=.5\ht\pgfnodeparttextbox % height of the box
    \pgf@x=.5\wd\pgfnodeparttextbox % width of the box
    \addtolength{\pgf@x}{2ex}
    \setlength{\pgf@xa}{\pgfshapeminwidth}
    \ifdim\pgf@x<.5\pgf@xa
    \pgf@x=.5\pgf@xa
    \fi
    \setlength{\pgf@ya}{\pgfshapeminheight}
    \ifdim\pgf@y<.5\pgf@ya
    \pgf@y=.5\pgf@ya
    \fi
}
\savedanchor{\southwest}{
    \pgf@y=-.5\ht\pgfnodeparttextbox
    \pgf@x=-.5\wd\pgfnodeparttextbox
    \addtolength{\pgf@x}{-2ex}
%     \pgf@x=0pt
    \setlength{\pgf@xa}{\pgfshapeminwidth}
    \ifdim\pgf@x>-.5\pgf@xa
    \pgf@x=-.5\pgf@xa
    \fi
    \setlength{\pgf@ya}{\pgfshapeminheight}
    \ifdim\pgf@y>-.5\pgf@ya
    \pgf@y=-.5\pgf@ya
    \fi
}
\anchor{text}{%
    \northeast%
    \pgf@x=-.5\wd\pgfnodeparttextbox%
    \pgfmathsetlength{\pgf@y}{-.5ex}
}
\anchor{north east}{
    \northeast
    \pgfmathsetmacro{\nw}{(1-sin(45))*\neoffset}
    \addtolength{\pgf@x}{-\nw pt}
    \addtolength{\pgf@y}{-\nw pt}
}
\anchor{center}{
    \pgf@x=0pt
    \pgf@y=0pt
}
\anchor{south west}{
    \southwest
    \pgfmathsetmacro{\nw}{(1-sin(45))*\swoffset}
    \addtolength{\pgf@x}{\nw pt}
    \addtolength{\pgf@y}{\nw pt}
}
\anchor{north west}{
    \northeast
    \pgfmathsetmacro{\temp@x}{\pgf@x}
    \southwest
    \pgfmathsetmacro{\temp@xtwo}{\pgf@x}
    \northeast
    \pgfmathsetmacro{\xdiff}{\temp@x-\temp@xtwo}
    \def\pgf@xa{\pgf@x-\xdiff}
            \pgfmathsetmacro{\nw}{(1-sin(45))*\nwoffset}
    \def\pgf@xaa{\pgf@xa+\nw}
    \def\pgf@yaa{\pgf@y-\nw}
    \pgfpoint{\pgf@xaa}{\pgf@yaa}
}
\anchor{south east}{
    \southwest
    \pgfmathsetmacro{\temp@x}{\pgf@x}
    \northeast
    \pgfmathsetmacro{\temp@xtwo}{\pgf@x}
    \southwest
    \pgfmathsetmacro{\xdiff}{\temp@x-\temp@xtwo}
    \def\pgf@xa{\pgf@x-\xdiff}
    \pgfmathsetmacro{\nw}{(1-sin(45))*\seoffset}
    \def\pgf@xaa{\pgf@xa-\nw}
    \def\pgf@yaa{\pgf@y+\nw}
    \pgfpoint{\pgf@xaa}{\pgf@yaa}
}
\anchor{south}{\south}
\anchor{north}{\north}
\anchor{east}{\east}
\anchor{west}{\west}
\backgroundpath{% this is new
    % store lower right in xa/ya and upper right in xb/yb
    \southwest \pgf@xa=\pgf@x \pgf@ya=\pgf@y
    \northeast \pgf@xb=\pgf@x \pgf@yb=\pgf@y
    % construct main path
    \pgfkeysgetvalue{/tikz/rectangle with rounded corners north west}{\pgf@rectc}
    \pgfsetcornersarced{\pgfpoint{\pgf@rectc}{\pgf@rectc}}
    \pgfpathmoveto{\pgfpoint{\pgf@xa}{\pgf@ya}}
    \pgfpathlineto{\pgfpoint{\pgf@xa}{\pgf@yb}}
    \pgfkeysgetvalue{/tikz/rectangle with rounded corners north east}{\pgf@rectc}
    \pgfsetcornersarced{\pgfpoint{\pgf@rectc}{\pgf@rectc}}
    \pgfpathlineto{\pgfpoint{\pgf@xb}{\pgf@yb}}
    \pgfkeysgetvalue{/tikz/rectangle with rounded corners south east}{\pgf@rectc}
    \pgfsetcornersarced{\pgfpoint{\pgf@rectc}{\pgf@rectc}}
    \pgfpathlineto{\pgfpoint{\pgf@xb}{\pgf@ya}}
    \pgfkeysgetvalue{/tikz/rectangle with rounded corners south west}{\pgf@rectc}
    \pgfsetcornersarced{\pgfpoint{\pgf@rectc}{\pgf@rectc}}
    \pgfpathclose
}
}
\makeatother